% Generated draft: requires mathematical and source review.
\documentclass[11pt]{article}
\usepackage[margin=25mm]{geometry}
\usepackage{fontspec}
\setmainfont{Latin Modern Roman}
\usepackage{amsmath,amssymb,mathtools,unicode-math}
\setmathfont{Latin Modern Math}
\usepackage{enumitem,needspace,xurl,hyperref}
\setlength{\parindent}{0pt}
\setlength{\parskip}{5pt}
\setlength{\emergencystretch}{3em}
\newcommand{\N}{\mathbb{N}}
\newcommand{\Z}{\mathbb{Z}}
\newcommand{\Q}{\mathbb{Q}}
\newcommand{\R}{\mathbb{R}}
\newcommand{\C}{\mathbb{C}}
\newcommand{\F}{\mathbb{F}}
\newcommand{\A}{\mathbb{A}}
\newcommand{\PP}{\mathbb{P}}
\newcommand{\E}{\mathbb{E}}
\newcommand{\eps}{\varepsilon}
\newcommand{\dd}{\,\mathrm{d}}
\begin{document}
\section*{UnsolvedMath ID: 30002490}
\textbf{Source problem code:} catalogue entry
\subsection{Definitions and mathematical statement}
The supplied record poses a research programme rather than a fully specified conjecture. Its subject is finite-time invariant manifolds in aperiodic, unsteady flows and the transport associated with stable and unstable manifold segments.

For notation only, let \(I\subseteq\R\) denote a time interval and let \(M\) denote an unspecified state space. Write \(\Phi_{t,s}\) for evolution from time \(s\) to time \(t\), wherever that evolution is defined. The usual evolution identities are
\[
\Phi_{s,s}=\operatorname{id},\qquad
\Phi_{t,r}\circ\Phi_{r,s}=\Phi_{t,s},
\]
wherever the compositions exist. A family of sets \(W(t)\subseteq M\) is transported by this evolution if
\[
\Phi_{t,s}\bigl(W(s)\bigr)=W(t)
\]
for the relevant times and domains. This identity supplies notation for material evolution; it does not by itself define finite-time stability, instability, or a transport barrier. Neither \(M\), the admissible evolutions, nor the observation interval is specified by the supplied record.

\textbf{Source-qualified research question.} Develop a theory of finite-time invariant manifolds that captures the time dependence of aperiodic unsteady flows. For stable and unstable manifold segments, denoted provisionally by \(W^s(t)\) and \(W^u(t)\), quantify transport associated with their intersections or their absence, and relate these structures to practical diagnostics of flow barriers. In particular, address the source fragment's question: does the proposed description capture time variation, and how do intersections or their absence produce transport?

The superscripts \(s\) and \(u\) are labels for the structures requested in the record, not definitions of finite-time stability. The supplied fragment does not identify the proposed description. It also does not specify stability criteria, endpoint conditions, a transported quantity or measure, regions between which transport is measured, or a criterion for agreement with barrier diagnostics. Consequently, it does not determine a precise universally or existentially quantified theorem. These choices remain unresolved parts of the formulation; no particular choice is asserted here.
\subsection{Short English statement}
The problem asks for a finite-time theory that describes changing structures in aperiodic unsteady flows, quantifies transport associated with intersecting or nonintersecting stable and unstable manifold segments, and connects those structures to practical flow-barrier diagnostics. The supplied fragment leaves essential definitions and hypotheses unspecified, so this is a source-qualified research programme requiring mathematical review.
\subsection{Sources}
\begin{itemize}
\item Supplied local UnsolvedMath registry record: permanent ID \(30002490\), exact source problem code \texttt{catalogue entry}, title \emph{Finite-Time Invariant Manifolds and Unsteady Transport}, in the Oberwolfach Reports Open Problems collection. The record assigns the year 2014 and supplies no individual proposer.
\item Supplied report citation: \emph{Mixing, Transport and Coherent Structures} (2014), Oberwolfach Reports, pp.\ 213--286, DOI \texttt{10.4171/owr/2014/04}. Supplied source URL and literature-source link: \url{https://doi.org/10.4171/owr/2014/04}. The background describes Oberwolfach Reports as documenting mathematical workshops and their open-problem discussions.
\item Additional supplied literature source, labelled ``Lagrangian coherent-structure review'': \url{https://doi.org/10.1146/annurev-fluid-010313-141322}. No authors or fuller bibliographic details are supplied for this reference.
\item Literature assessment dated exactly \texttt{2026-08-21}: the supplied curation labels the problem open, states that the report identifies the questions as unresolved, and reports that no general resolution was found. This is supplied dated curation, not an independently verified literature assessment. The background explicitly requests independent verification of time-sensitive status and cited claims.
\item The registry labels the statement \texttt{corrected\_verified} and reports that it was checked against the official Oberwolfach report, with the clean statement restoring source hypotheses and notation. These are supplied verification claims; no independent source inspection or mathematical verification is asserted in this draft.
\end{itemize}
\end{document}
